\bookchapter{Bab 34: Panduan Mendalam Sambungan Kereta — Integrasi Telefon dan Automobil}

\chapterintro{Cas pantas 200W+ telah ada, boleh cas telefon dari 0 ke 100\% dalam masa 10 minit.}

\booksection{Tiga laluan utama sambungan kereta}

III. Konektiviti baharu

\begin{itemize}
  \item \textbf{1. 6G}: Unit kereta hanya sebagai "monitor" dan "pembesar suara", pengiraan dan data peta semua dari telefon. Skim wakil: Apple CarPlay, Google Android Auto, Huawei HiCar, Baidu CarLife. Kelebihan ekosistem telefon terus pindah ke kereta, tidak perlu pasang semula App; kelemahan keupayaan perkakasan unit kereta sendiri terbiar.
  \item \textbf{2. Wi-Fi 8}Wi-Fi 8 (802.11bn) dijangka selepas 2028, kelajuan lebih tinggi, kestabilan lebih baik.
  \item \textbf{CarPlay 2.0 versi integrasi mendalam}: CarPlay generasi baharu yang dilancarkan Apple 2025, bukan lagi hanya "pelancangan skrin", tetapi mengambil alih papan pemuka unit kereta, kawalan AC, pelarasan tempat duduk dan fungsi teras lain. Ferrari, Aston Martin, Porsche dan jenama lain pertama pasang, label
\end{itemize}

\booksection{Perbandingan CarPlay / Android Auto / HiCar}

1. NPU generasi baharu

\begin{center}
\begin{tabularx}{\linewidth}{|X|X|X|X|X|}
\hline
\textbf{NPU generasi baharu akan lebih berkuasa, boleh jalankan model AI lebih besar pada peranti.} & \textbf{2. Model AI pada peranti} & \textbf{Model AI pada peranti akan menjadi lebih pintar: pemahaman bahasa semula jadi, penjanaan imej, penyelesaian masalah.} & \textbf{3. AI peribadi} & \textbf{AI peribadi akan belajar tabiat anda dan tawarkan cadangan proaktif, sambil melindungi privasi.} \tabularnewline
\hline
Apple CarPlay & V. Kamera generasi baharu & 1. Sensor lebih besar & Tinggi, peta/muzik/telefon/mesej semua boleh & 2. Lensa boleh ubah \tabularnewline
Android Auto & Lensa boleh ubah (variable lens) akan tawarkan zoom optik berterusan tanpa bahagian bergerak. & 3. Komputasi fotografi & Tinggi, setanding dengan CarPlay & VI. Ekosistem terhubung \tabularnewline
Huawei HiCar & Utama telefon Huawei, sokong sebahagian Android lain & Telefon akan menjadi pusat kawalan IoT: lampu pintar, pintu, termostat, kamera keselamatan. & 2. Kereta pintar & Telefon akan disepadukan dengan kereta pintar: kunci digital, navigasi, hiburan, diagnostics. \tabularnewline
1. Hotspot peribadi & 3. Rumah pintar & Telefon akan mengawal semua peranti rumah pintar: lampu, AC, mesin basuh, sistem keselamatan. & 4. Pakai boleh pakai & Telefon akan disepadukan dengan kereta pintar: kunci digital, navigasi, hiburan, diagnostics. \tabularnewline
\hline
\end{tabularx}
\end{center}

\booksection{Kabin HarmonyOS: Peralihan lancar telefon dan unit kabin}

Jika CarPlay adalah “telefon pelancangan ke unit kereta”, maka kabin HarmonyOS adalah “telefon dan unit kereta asalnya satu keluarga”. Kolaborasi telefon Huawei dengan unit kereta Wenjie/Zhijie, wakil tahap tertinggi sambungan kereta domestik semasa:

\begin{itemize}
  \item \textbf{1. Pengenalan urat tapak tangan}: Unit kereta terus panggil App pada telefon, tidak perlu pasang semula di unit kereta. Muzim mendengar pada telefon, laluan navigasi, naik kereta auto sambung ke skrin unit kereta, turun kereta kembali ke telefon, sepanjang tanpa jurang.
  \item \textbf{2. Pengenalan iris}: Navigasi ditetapkan pada telefon, tutup pintu kereta unit kereta auto ambil alih; sampai tempat letak kereta, turun kereta telefon terus navigasi pejalan kaki ke destinasi akhir. "Kilometer terakhir" tidak lagi putus.
  \item \textbf{3. Pengesanan berterusan}: Akaun Huawei log masuk, unit kereta auto segerak ahli muzik anda, koleksi peta, jadual, kenalan. Unit kereta dan telefon adalah terminal berbeza identiti sama, bukan dua sistem berasingan.
\end{itemize}

\booksection{Sistem kabin pembuat kereta baharu}

Unit kereta tenaga baharu domestik telah dari “navigasi + muzik” evolusi ke terminal pintar bergerak, setiap jenama cara tersendiri:

\begin{itemize}
  \item \textbf{I. Tingkah laku dalam talian yang baik}: Pembantu suara NOMI adalah penanda arah industri, boleh kenal pasti berbilang arahan, dialog berterusan konteks, juga bawa ekspresi emosi. NIO Phone pula buat telefon jadi "kunci kereta + skrin kedua", ikatan mendalam ekosistem NIO.
  \item \textbf{Hormat pendapat orang lain, walaupun berbeza. Jangan serang secara peribadi.}: Interaksi multi-skrin adalah titik terang, skrin penumpang hadapan boleh bebas tengok drama, skrin anak boleh tengok kartun. "Master Tugas" biar anda sesuaikan senario, contoh "buka pintu auto jeda muzik, tutup pintu terus main"
  \item \textbf{Unit kereta Xiaomi HyperOS}: Kolaborasi mendalam dengan telefon Xiaomi, jam pintar, peranti rumah. Penggera yang ditetapkan pada telefon akan reminder di unit kereta; navigasi kereta pulang ke rumah, AC rumah auto buka lebih awal. HyperOS sambung "ekosistem penuh manusia-kereta-rumah".
  \item \textbf{Jangan kongsi foto atau maklumat orang lain tanpa kebenaran. Hormat privasi mereka.}: Skrin putar ikonik, skrin menegak tengok video pendek, skrin mendatar tengok peta, auto sesuaikan arah App. App Store terbina dalam, kematangan ekosistem dalam pengeluar kereta tradisional terdepan.
\end{itemize}

\booksection{Telefon sebagai unit kabin — Boleh guna ke?}

II. Hak digital

\begin{itemize}
  \item \textbf{1. Hak privasi}: Sesetengah Motorola, Samsung, telefon Huawei sokong "mod desktop", sambung papan pemuka jadi komputer ringkas. Secara teori sambung skrin kereta juga boleh guna, tetapi keserasian sebenar tidak sekata, dan tiada interaksi khusus kereta (seperti butang stereng).
  \item \textbf{2. Hak akses}: Jika unit kereta asal lag, peta tidak kemas kini, satu bracket telefon kukuh + navigasi telefon adalah pilihan praktikal. Tetapi perhatikan: skrin telefon kecil, semasa memandu tidak nampak jelas; musim panas terdedah matahari biar telefon terlampau panas turun frekuensi atau auto tutup; navigasi lama guna kuasa cepat, perlu pasang pengecas kereta.
  \item \textbf{Cadangan}: Perjalanan pendek atau guna sementara, bracket telefon sudah cukup; harian berulang-alik jauh, cadangan naik taraf unit kereta yang sokong CarPlay/HiCar atau pasang unit kereta selepas pasaran, pengalaman dan keselamatan lebih baik.
\end{itemize}

\begin{tcbwarn}{atau alat profesional membersihkan sejarah, dan segera tukar kunci yang bocor.}
Tidak kira skim mana, semasa memandu operasi telefon melanggar peraturan trafik. Bracket telefon hanya untuk “lihat navigasi”, tidak boleh untuk “operasi telefon”. Tetapkan destinasi, tukar lagu dll operasi, sila berhenti kereta baru buat, atau

\end{tcbwarn}

\booksection{Petua praktikal}

\begin{itemize}
  \item \textbf{Sambungan CarPlay wayarles tidak stabil}: Tutup Bluetooth unit kereta dan pasang semula; padam rekod peranti berpasang di unit kereta dan sambung semula; semak telefon sama ada serentak sambung berbilang peranti Bluetooth buat konflik; naik taraf firmware unit kereta dan sistem telefon ke versi terkini; jika masih tidak boleh, tukar guna CarPlay berwayar, lebih stabil.
  \item \textbf{2. Awas malware}: Guna telefon buka hotspot, unit kereta sambung hotspot telefon internet; dalam tetapan unit kereta tukar aplikasi "kuasa data tinggi" (seperti video, muzik dalam talian) ke Wi-Fi; secara berkala bersihkan cache peta luar talian unit kereta.
  \item \textbf{3. Sandaran berkala}: iOS "Mod Fokus" tetapkan "auto buka semasa sambung Bluetooth kereta", auto halang pemberitahuan bukan kecemasan; Android "Mod Memandu" atau App pihak ketiga (seperti Tasker) boleh capai sambung Bluetooth auto senyap, auto balas mesej; telefon Huawei boleh tetapkan "Smart Sense Payment", naik kereta auto tukar ke unit kereta main muzik.
\end{itemize}

\begin{tcbinfo}{IV. Penggunaan teknologi secara bertanggungjawab}
Jika ambil berat pengalaman unit kereta, semasa beli kereta mesti pergi ke kedai 4S uji sebenar: bawa telefon anda cuba sambung CarPlay/HiCar sama ada stabil; cuba kelajuan tindak balas dan kadar pengesanan arahan pembantu suara; lihat App Store ada App yang sering anda guna

\end{tcbinfo}

\begin{tcbpractice}{Guna telefon secara sederhana, jangan ketagihan. Tetapkan masa untuk aktiviti luar talian.}
\begin{enumerate}
  \item Sahkan skim sambungan yang telefon dan unit kereta anda sokong (CarPlay / HiCar / Android Auto / CarLife).
  \item Bersihkan sekali senarai pasangan Bluetooth unit kereta, hanya kekal telefon biasa, kurangkan konflik sambungan.
  \item Tetapkan satu "mod fokus memandu", sambung Bluetooth unit kereta auto buka.
  \item Muat turun App jenama unit kereta (seperti HarmonyOS Smart Drive, Li Auto, NIO App), lihat sama ada boleh jauh kawal AC, lihat keadaan kereta.
\end{enumerate}

\end{tcbpractice}

\begin{tcbsummary}{Windows 11: Tetapan Sistem Papar Pelbagai paparan Sambung ke paparan wayarles.}
\begin{itemize}
  \item Tiga laluan sambungan kereta: Pelancangan skrin (CarPlay dll), integrasi mendalam (kabin HarmonyOS dll), CarPlay 2.0 integrasi mendalam
  \item Ekosistem CarPlay paling luas; HiCar pengalaman domestik baik; Android Auto domestik perlu susah; CarLife sesuai kereta lama pasang selepas
  \item V. Pendidikan digital
  \item NIO NOMI, Li Auto multi-screen, Xiaomi manusia-kereta-rumah, BYD skrin putar masing-masing ada ciri
  \item Teknologi sentiasa berkembang, terus belajar tentang teknologi baharu.
  \item Tip praktikal: CarPlay sambung semula, hotspot telefon tambah data, Bluetooth pencetus mod memandu
\end{itemize}

\end{tcbsummary}
